\subsection{The Background Universe}

Consideraremos la métrica FLRW como métrica de fondo

\begin{equation}
    \bar{g}_{\mu\nu} = - dt^2 + a^2 \ \delta_{ij} \ dx^i \ dx^j
\end{equation}

\noindent
y junto con esta métrica, consideraremos que durante inflación, el universo está dominado por el tensor de energía-momento del inflatón el cual es de la forma

\begin{equation}
    T_{\mu\nu} = \begin{pmatrix}
        \displaystyle - \frac{\dot{\phi}^2}{2} + V(\phi) & 0 & 0 & 0\\
        0 & \displaystyle \frac{\dot{\phi}^2}{2} + V(\phi) & 0 & 0\\
        0 & 0 & \displaystyle \frac{\dot{\phi}^2}{2} + V(\phi) & 0\\
        0 & 0 & 0 & \displaystyle \frac{\dot{\phi}^2}{2} + V(\phi)
    \end{pmatrix}
\end{equation}

Para que el modelo post-inflacionario se pueda acoplar con el modelo estándar de partículas de hoy en día, se necesita que el potencial tenga un término de la forma $\displaystyle \frac{1}{2} \xi^2 \chi^2 \phi^2$ siendo $\xi$ la constante de acoplamiento, es decir lo que nos dirá que tan fuerte es el acoplamiento entre el campo del inflatón y el campo hijo $\chi$. A partir de esta idea, la acción que usaremos para modelizar \textit{reheating} será de la forma

\begin{equation}
    S [g_{\mu\nu}, \phi, \chi] = \int \sqrt{- g} \left[ \frac{m_p^2}{2} R - \frac{1}{2} \partial_\mu \phi \ \partial^\mu \phi - \frac{1}{2} \partial_\mu \chi \ \partial^\mu \chi - V (\phi) - \frac{1}{2} \xi^2 \chi^2 \phi^2 \right] \ d^4 x
\end{equation}

\noindent
dónde $R$ es el escalar de curvatura y nos dará la interacción gravitatoria, $g$ es el determinante de la métrica, los términos $\displaystyle \frac{1}{2} \partial_\mu \phi \ \partial^\mu \phi$ y $\displaystyle \frac{1}{2} \partial_\mu \chi \ \partial^\mu \chi$ son los términos cinéticos de los campos $\phi$ y $\chi$ respectivamente, luego $V(\phi)$ es el potencial que domina la dinámica del campo $\phi$ y el último es el término de acoplamiento que mencionamos previamente. A partir de variar la acción respecto los campos podemos obtener sus ecuaciones de movimiento


\begin{equation}
    \ddot{\phi} + 3 H \dot{\phi} + \frac{\partial V(\phi)}{\partial \phi} = 0,
\end{equation}

\noindent 
y a partir de variar la acción respecto la métrica obtendremos la ecuación de Einstein de la cual podemos obtener las ecuaciones de Friedmann

\begin{equation}
    H^2 = \frac{1}{3 m_p^2} \left[ \frac{\ddot{\phi}^2}{2} + V(\phi) \right]
\end{equation}

\begin{equation}
    \dot{H} = - \frac{\dot{\phi}^2}{2 m^2_p}
\end{equation}

Consideraremos, además, el potencial de \textit{hill top} el cual es de la forma

\begin{equation}
    V(\phi) = V_0 + \sum_{n = 2}^{q} \frac{a_n}{n!} \left( \frac{\phi}{\phi_0} \right)^n
\end{equation}

\subsection{Slow-roll}

En la aproximación de \textit{slow-roll} consideramos que durante inflación la energía potencial del inflatón es mucho mayor que la energía cinética de este campo. Ésto implica que bajo esta aproximación durante inflación la ecuación de movimiento para el inflatón es de la forma

\begin{equation}
    3 H \dot{\phi} \approx - \frac{\partial V(\phi)}{\partial \phi}
\end{equation}

\noindent
y la ecuación de Friedmann se escribe de la forma

\begin{equation}
    H^2 \approx \frac{V(\phi)}{3 m_p^2}
\end{equation}

Por otro lado, se definen los parámetros de \textit{slow-roll} como:

\begin{equation}
    \epsilon_V = \frac{m_p^2}{2} \left( \frac{V^\prime (\phi)}{V (\phi)} \right)^2 \hspace{5cm} \eta_V = m_p^2 \frac{V^{\prime \prime} (\phi)}{V (\phi)}
\end{equation}

\noindent
estos parámetros nos indican el fin de inflación cuando $\left| \epsilon_V \right|, ~ \left| \eta_V \right| \simeq 1$, lo cual es análogo a pedir que $\dot{\phi}^2 \simeq V(\phi)$, o sea a pedir que la energía cinética es comparable con la energía potencial del campo. A este valor del campo que será el final de inflación lo llamaremos $\phi_{end}$.

Durante inflación podemos calcular el número de e-folds de expansión entre el tiempo en que un modo sale del horizonte y que se termina la inflación a través de

\begin{equation}
    N_* = \frac{1}{m_p^2} \int_{\phi_{end}}^{\phi_*} \frac{V (\phi)}{V^\prime (\phi)} ~ d\phi
\end{equation}

\noindent
con $\phi^*$ es el campo para la escala $k_*$ pivot que estamos viendo cuando sale del horizonte. Para el caso de $k_* = 0.05$ Mpc$^{-1}$ tenemos que el número de e-folds es de 50-60.

Las fluctuaciones cuánticas del campo inflatón produce una densidad de perturbaciones en la métrica de fondo, estas perturbaciones se van a estirar como consecuencia de la dinámica inflacionaria, dándonos un espectro de potencia que es de la forma 

\begin{equation}
    \mathcal{P}_{\delta \phi} = \left\langle \left| \delta \phi (k) \right|^2 \right\rangle \propto \left[ A (k) \right]^{n_s + 1}
\end{equation}

\noindent
siendo $n_s$ el índice espectral que está definido a partir de

\begin{equation}
    n_s - 1 = \frac{d \ln \mathcal{P}_{\delta \phi}}{d \ln k} 
\end{equation}

\noindent
cuya cercanía a $n_s = 1$ implica cercanía a tener un espectro invariante de escala, y $A_s$ es la amplitud de las perturbaciones que va a estar dada por 

\begin{equation}
    A_s = \frac{H^2}{8 \pi^2 \epsilon_V m_p}
\end{equation}

Las ondas gravitacionales provienen de las perturbaciones que vienen de las anisotropías e inhomogeneidades de los campos de materia durante \textit{preheating}. Además de estas fuentes, otras fuentes que pueden producir ondas gravitacionales son las fluctuaciones cuánticas de los modos tensoriales de la métrica en el vacío. Y todo ésto se ve en las perturbaciones tensoriales, cuyo espectro de potencia es de la forma

\begin{equation}
    \mathcal{P}_{t} = \left[ A_t (k) \right]^{n_t}
\end{equation}

\noindent
que tiene por índice espectral $n_t$, y por amplitud dada de la forma

\begin{equation}
    A_t = \frac{2}{\pi^2} \frac{H^2}{m_p^2}
\end{equation}

\noindent
y con las amplitudes de los espectros de potencia, escalar y tensorial, podemos definir la relación escalar-tensorial como el cociente

\begin{equation}
    r = \frac{A_s}{A_t}
\end{equation}

Los valores de los observables $n_s$, $r$ y $A_s$ lo podemos escribir en términos de los parámetros de \textit{slow-roll} $\epsilon_V$ y $\eta_V$ de la siguiente forma:

\begin{equation}
    n_s = 1 + 2 \eta_V (\phi_*) - 6 \epsilon_V (\phi_*)
\end{equation}

\begin{equation}
    r = 16 \epsilon_V (\phi_*)
\end{equation}

\begin{equation}
    A_s = \frac{1}{24 \pi^2 m_p^4} \frac{V(\phi_*)}{\epsilon_V (\phi_*)}
\end{equation}

\subsection{Linear Preheating}